\documentclass{article}
\usepackage[left=1in, right=1in, bottom = 1in, top = 1in]{geometry}
\usepackage{graphicx} % Required for inserting images
\usepackage{bbm}
\usepackage{amsmath}
\usepackage[english]{babel}
\usepackage[utf8]{inputenc}
\usepackage[authoryear]{natbib}
%\usepackage[margin=1in]{geometry}
\usepackage{setspace}

%\setlength{\parskip}{1em}
%\onehalfspacing
%\linespread{1.45}

% Math Packages
\usepackage{amsmath}
%\numberwithin{equation}{section}
\usepackage{amsfonts}
\usepackage{amsthm}
\usepackage{amssymb}
\usepackage{mathrsfs}
\usepackage{mathtools}
%\usepackage{bbm}
\usepackage{amstext}
\usepackage{comment}
\usepackage{epsfig}
\usepackage{latexsym}
\usepackage{epstopdf}
\usepackage{bm}
\usepackage{algorithm}
\usepackage{algpseudocode}
%\usepackage[capposition=top]{floatrow}


% For drawing graph
\usepackage{tikz}
\usetikzlibrary{positioning,chains,fit,shapes,calc}
\definecolor{myblue}{RGB}{80,80,160}
\definecolor{mygreen}{RGB}{80,160,80}

\usepackage{url}
\usepackage[colorlinks = true, linkcolor = blue]{hyperref}
\hypersetup{
        urlcolor = blue,
	colorlinks = true,
	linkcolor = blue,
	citecolor = blue,
}
\usepackage{paralist}
\usepackage{enumitem}
\setlist[enumerate]{topsep=0pt,noitemsep}

% For Tables
\usepackage{multirow}
\usepackage{multicol}
\usepackage{babel}
\usepackage{array}
\usepackage{rotating}
\usepackage{longtable}
\usepackage{float}
\usepackage{booktabs}
\usepackage{lscape}
\usepackage{caption}
\usepackage{subcaption}
\usepackage{afterpage}
% \usepackage{showkeys}
\usepackage{authblk}

% Theorem Environment
\newtheorem{theorem}{Theorem}[section]
\newtheorem{remark}{Remark}[section]
\newtheorem{assumption}{Assumption}
\newtheorem{lemma}{Lemma}[section]
\newtheorem{notation}{Notation}
\newtheorem{example}{Example}
\newtheorem{definition}{Definition}[section]
\newtheorem{proposition}{Proposition}[section]
\newtheorem{corollary}{Corollary}[section]
\newtheorem{prediction}{Prediction}
\newtheorem{regression}{Regression}

% Custom environment for continuance
\newenvironment{continuance}[1]
{%
  \newcommand{\continuanceref}{\ref{#1}}% Define reference to original example
  \continuancex % Begin the non-numbered theorem-like environment
}
{\endcontinuancex}

% Define a non-numbered continuation environment
\newtheorem*{continuancex}{Example \continuanceref \ (continued)}

% \makeatletter
% \renewenvironment{proof}[1][\proofname]{%
%   \par\pushQED{\qed}\normalfont%
%   \topsep6\p@\@plus6\p@\relax
%   \trivlist\item[\hskip\labelsep\bfseries#1\@addpunct{.}]%
%   \ignorespaces
% }{%
%   \popQED\endtrivlist\@endpefalse
% }
% \makeatother
%\numberwithin{theorem}{section}
%\numberwithin{corollary}{section}
%\numberwithin{proposition}{section}
%\numberwithin{lemma}{section}
%\numberwithin{assumption}{section}

% Formatting 
\linespread{1.4}
% \mathtoolsset{showonlyrefs=true}

\allowdisplaybreaks


% Define math operators
\DeclareMathOperator*{\argmax}{argmax}
\DeclareMathOperator*{\argmin}{argmin}
\DeclareMathOperator{\tr}{tr}
\DeclareMathOperator{\vect}{vec}
\DeclareMathOperator{\Rank}{Rank}

% Define shortcut
\newcommand{\E}{\mathbb{E}}
\newcommand{\B}{B}
\newcommand{\Var}{\mathrm{Var}}
\newcommand{\Cov}{\mathrm{Cov}}
\newcommand{\bbm}[1]{\mathbbm{#1}}
\newcommand{\bb}[1]{\mathbb{#1}}
\newcommand{\floor}[1]{\lfloor #1 \rfloor}
\newcommand{\sq}[1]{\textquoteleft #1 \textquoteright}
\newcommand{\dq}[1]{\textquotedblleft #1 \textquotedblright}
\newcommand{\red}{\color{red}}
\newcommand{\norm}[1]{\left\lVert#1\right\rVert}
\def\T{{\textsc{t}}}

\newcommand{\pr}{\textup{pr}}
\newcommand{\cov}{\textup{cov}}
\newcommand{\var}{\textup{var}}
\newcommand{\avar}{\textup{avar}}
\newcommand{\acov}{\textup{acov}}
\newcommand{\kprime}{k^{\prime}}
\newcommand{\adj}{\textup{adj}}

\newcommand{\bOneZ}{\boldsymbol{1}_{m}}
\newcommand{\bZeroZ}{\boldsymbol{0}_{m}}
\newcommand{\bOneA}{\boldsymbol{1}_{n}}
\newcommand{\bZeroA}{\boldsymbol{0}_{n}}
\newcommand{\Yt}[1]{Y_{#1}(\boldsymbol{1}_{m})}
\newcommand{\Yc}[1]{Y_{#1}(\boldsymbol{0}_{m})}

\newcommand{\bY}{\boldsymbol{Y}}
\newcommand{\bz}{\boldsymbol{z}}
\newcommand{\bZ}{\boldsymbol{Z}}
\newcommand{\ba}{\boldsymbol{a}}
\newcommand{\bA}{\boldsymbol{A}}
\newcommand{\bW}{\boldsymbol{W}}
\newcommand{\bV}{\boldsymbol{V}}
\newcommand{\bX}{\boldsymbol{X}}
\newcommand{\bpi}{\boldsymbol{\pi}}
\newcommand{\Davg}{D_{\textup{avg}}}
\newcommand{\bOmega}{\boldsymbol{\Omega}}
\newcommand{\bLambda}{\boldsymbol{\Lambda}}

\newcommand{\cN}{\mathcal{N}}
\newcommand{\cF}{\mathscr{F}}
\newcommand{\cT}{\mathcal{T}}
\newcommand{\cS}{\mathcal{S}}

\newcommand{\tA}{\tilde{A}}
\newcommand{\tS}{\tilde{S}}
\newcommand{\tH}{\tilde{H}}


\newcommand{\allt}{\left(\prod_{k\in \mathcal{G}_{+i}} Z_k\right)}
\newcommand{\alltj}{\left(\prod_{k\in \mathcal{G}_{+j}} Z_k\right)}
\newcommand{\allc}{\left\{\prod_{k\in \mathcal{G}_{+i}} (1-Z_k)\right\}}
\newcommand{\allcj}{\left\{\prod_{k\in \mathcal{G}_{+j}} (1-Z_k)\right\}}

\newcommand{\plim}{\textup{plim}}

\newcommand{\bG}{\boldsymbol{G}}
\newcommand{\bGamma}{\boldsymbol{\Gamma}}
\newcommand{\G}{G}
\newcommand{\cG}{\mathcal{G}}
\newcommand{\cGijcap}{\cG_{+i} \cap \cG_{+j}}
\newcommand{\cGijcup}{\cG_{+i} \cup \cG_{+j}}

\newcommand{\Smax}{{\max_{1\leq i\leq n}\G_{+i}}}
\newcommand{\Dmax}{{\max_{1\leq k\leq m}\G_{k+}}}
\newcommand{\Smaxsupp}{{\bar{D}_\textup{O}}}
\newcommand{\Dmaxsupp}{{\bar{D}_\textup{I}}}

\newcommand{\sumi}{\sum_{i=1}^n}
\newcommand{\sumj}{\sum_{j=1}^n}
\newcommand{\sumk}{\sum_{k=1}^m}
\newcommand{\sumij}{\sum_{i=1}^{n} \sum_{j=1}^{n}}
\newcommand{\sumjneqi}{\sum_{j\neq i}}
\newcommand{\sumu}{\sum_{u=1}^n}
\newcommand{\sumv}{\sum_{v=1}^n}
\newcommand{\maxi}{\max_{1\leq i\leq n}}

\newcommand{\de}{\textsc{de}}
\newcommand{\bie}{\textsc{bie}}
\newcommand{\wie}{\textsc{wie}}
\newcommand{\horthom}{\textup{ht}}
\newcommand{\haj}{\textup{haj}}
\newcommand{\diag}{\textup{diag}}

\newcommand{\indicator}{\mathbbm{1}}
\newcommand{\bbI}{\mathbbm{I}}
\newcommand{\ashprime}{a^{\prime},s^{\prime},h^{\prime}}

% propensity scoressssss
\newcommand{\pscore}[1]{\pi_{#1}}

\newcommand{\sz}[1]{{\textcolor{olive}{(Sizhu: #1)}}}

\makeatletter
\newcommand{\ostar}{\mathbin{\mathpalette\make@circled\star}}
\newcommand{\make@circled}[2]{%
  \ooalign{$\m@th#1\smallbigcirc{#1}$\cr\hidewidth$\m@th#1#2$\hidewidth\cr}%
}
\newcommand{\smallbigcirc}[1]{%
  \vcenter{\hbox{\scalebox{0.77778}{$\m@th#1\bigcirc$}}}%
}
\makeatother

\allowdisplaybreaks

%%%%%%%%%%%%%%%%%%%%%%%%%%%%%%%%%%%%%%%%%

\title{Estimating within-cluster and between-cluster spillover effects in randomized saturation designs}

% \author{ 
% Sizhu Lu\thanks{Department of Statistics, University of California, Berkeley, CA 94720, USA. Email: \href{mailto:sizhu_lu@berkeley.edu}{sizhu\_lu@berkeley.edu}}, 
% Lei Shi\thanks{Division of Biostatistics, University of California, Berkeley, CA 94720, USA. Email: \href{mailto:leishi@berkeley.edu}{leishi@berkeley.edu}}, 
% Peng Ding\thanks{Department of Statistics, University of California, Berkeley, CA 94720, USA. Email: \href{mailto:pengdingpku@berkeley.edu}{pengdingpku@berkeley.edu}}
% }

\author[1]{Sizhu Lu\thanks{The first two authors contributed equally to this work.}}
\author[2]{Lei Shi$^{*}$}
\author[1]{Peng Ding}

\affil[1]{Department of Statistics, University of California, Berkeley}
\affil[2]{Division of Biostatistics, University of California, Berkeley}

\date{\today}
\begin{document}

\maketitle
\begin{abstract}
    Randomized saturation designs are two-stage experiments: they first randomly assign treatment probabilities over the clusters and then randomly assign the treatment to the units within the clusters. The existing literature on randomized saturation designs focuses on estimating within-cluster spillover effects by assuming away between-cluster spillover effects. However, the units may interact across clusters in practice. A leading example is that some units are geographically close to each other, so spillover effects arise across clusters. Based on the potential outcomes framework, we formulate the causal inference problem of estimating within-cluster and between-cluster spillover effects in randomized saturation designs. We clarify the causal estimands and establish rigorous statistical theory for inference. We also apply our theory to analyze a recent randomized saturation design of cash transfer on household expenditure in Kenya.
\end{abstract}
% {\small\hspace{2em} \textbf{Key Words: } }


\section{Introduction}
Randomized saturation designs are widely used to study causal effects in the presence of network interference \citep{hudgens2008toward}. Most existing studies focus on within-cluster interference and assume no between-cluster spillover effects, which is often an unrealistic assumption in applied settings. A leading example is the General Equilibrium (GE) study by \citet{egger2022general}, which implemented a randomized experiment to evaluate the impact of a cash transfer program in rural Kenya. The study used a randomized saturation design: in the first stage, sublocations were randomly assigned to either high- or low-saturation groups. In the second stage, within each high-saturation sublocation, two-thirds of villages were randomly assigned to treatment, while in low-saturation sublocations, one-third of villages were treated. All eligible households in treated villages received transfers. In this setting, there are two plausible sources of interference: (i) within-sublocation interference, where villages in the same sublocation affect each other’s outcomes, and (ii) between-sublocation interference, where villages in different sublocations, but located near each other, may also exert influence.

To formalize these ideas, let $k(i)$ denote the sublocation to which village $i$ belongs, and let $\cG_{+i}$ represent the set of villages that are geographically close to village $i$. We define the exposure mapping as $d_i(\bA) = (A_{i}, S_{i}, H_{i})$, where $S_i$ and $H_i$ are summaries of within-sublocation exposure and between-sublocation exposure \citep{aronow2017estimating}:
\begin{eqnarray*}
    S_{i} = f_i\{(A_{j}: k(j)=k(i)\}, \quad 
    H_{i} = g_i\{A_j: k(j)\neq k(i), j\in\cG_{+i}\}. 
\end{eqnarray*}
For example, we can use the following exposure mapping:  
\begin{eqnarray*}
    S_{i} &=& \indicator\left\{\frac{\sum_{j\neq i}\indicator\{k(j)=k(i)\} A_{j}}{\sum_{j\neq i}\indicator\{k(j)=k(i)\}} > \frac{1}{2}\right\},  \\
    H_{i} &=& \indicator \left\{\frac{\sumj \indicator\{k(j)\neq k(i), j\in\cG_{+i}\}A_{j}}{\sumj \indicator\{k(j)\neq k(i), j\in\cG_{+i}\}} > \frac{1}{2}\right\}.
\end{eqnarray*}
Here, $S_{i}$ indicates whether more than half of the villages are getting treated within the same sublocation as village $i$, while $H_{i}$ measures whether more than half villages are treated among those connected to village $i$ but located in different sublocations. In other words, $S_{i}$ captures within-sublocation interference, since the potential outcome of village $i$ may depend not only on its own treatment, but also on the treatments of other villages in the same sublocation. On the other hand, $H_{i}$ captures between-sublocation interference, reflecting the possibility that nearby sublocations influence each other through spillover effects \citep{egger2022general}. 

\section{Causal parameters of interest}
\label{sec::parameters}
\subsection{Controlled causal effects}

\paragraph{Controlled direct effects.}
We first consider the direct effect of treatment. Define
\begin{eqnarray*}
    \de(s,h) &=& n^{-1}\sumi Y_{i}(1,s,h) - n^{-1}\sumi Y_{i}(0,s,h).
\end{eqnarray*}
The quantity $\de(s,h)$ represents the \textit{controlled direct effect} of treatment $A_{i}$ on the outcome, while holding the exposure variables $(S_{i},H_{i})$ fixed at values $(s,h)$.

\paragraph{Controlled indirect effects.}
We consider two types of indirect effects corresponding to the two sources of interference: within-cluster and between-cluster spillovers.
% comment on the different mechanisms of these two types of spillover effects.

For within-cluster indirect effects, we define the \textit{within-cluster controlled indirect effect} of $S_{i}$ on the outcome as
\begin{eqnarray*}
    \wie(s,s^{\prime},h) &=& n^{-1}\sumi Y_{i}(0,s,h) - n^{-1}\sumi Y_{i}(0,s^{\prime},h).
\end{eqnarray*}
The quantity $\wie(s,s^{\prime},h)$ captures the effect of changing the proportion of treated villages in the same sublocation (from $s$ to $s^{\prime}$) on the outcome of village $i$, while holding its own treatment fixed at zero and the between-cluster exposure $H_{i}$ at level $h$. 

For between-cluster indirect effects, similarly, we define the \textit{between-cluster controlled indirect effect} of $H_{i}$ on the outcome as
\begin{eqnarray*}
    \bie(s,h,h^{\prime}) &=& n^{-1}\sumi Y_{i}(0,s,h) - n^{-1}\sumi Y_{i}(0,s,h^{\prime}).
\end{eqnarray*}
The quantity $\bie(s,h,h^{\prime})$ captures the effect of changing the level of between-cluster exposure from $h$ to $h^{\prime}$, with the own treatment fixed at zero and the within-cluster exposure $S_{i}$ fixed at level $s$. 

\subsection{In-policy causal effects}
In practice, we can report the aggregated version of the causal parameters, by marginalizing over the implemented policy distributions. 
\paragraph{Marginal direct effect.}
We first consider the direct effect of treatment. Define
\begin{eqnarray*}
    \de &=& n^{-1}\sumi E_{\bA_{(-i)}}\left\{Y_{i}(1,S_{i},H_{i})\right\} - n^{-1}\sumi E_{\bA_{(-i)}}\left\{Y_{i}(0,S_{i},H_{i})\right\}, 
\end{eqnarray*}
where $E_{\bA_{(-i)}}(\cdot)$ denotes the expectation with respect to the distribution of $\bA_{(-i)}$, the vector of treatment assignments for all villages other than $i$. 

$\de$ represents the \textit{marginal direct effect} of treatment $A_{i}$, marginalizing over the distribution of treatments for all other villages. Therefore, by definition, $\de$ depends on the treatment assignment mechanism.


\paragraph{Marginal indirect effects.}
We again consider the two types of indirect effects corresponding to the two sources of interference: within-cluster and between-cluster spillovers.
% comment on the different mechanisms of these two types of spillover effects.

For within-cluster indirect effects, we define the \textit{marginal within-cluster indirect effect} of $S_{i}$ on the outcome as
\begin{eqnarray*}
    \wie(s,s^{\prime}) &=& n^{-1}\sumi E_{\bA_{(-i)}\mid S_{i} = s}\left\{Y_{i}(0,s,H_{i})\right\} - n^{-1}\sumi E_{\bA_{(-i)}\mid S_{i} = s^{\prime}}\left\{Y_{i}(0,s^{\prime},H_{i})\right\},
\end{eqnarray*}
where $E_{\bA_{(-i)}\mid S_{i} = s}(\cdot)$ denotes expectation with respect to the distribution of $\bA_{(-i)}$ conditional on $S_{i}=s$. The $\wie(s,s^{\prime})$ is a marginal version of this effect, marginalized over the conditional distribution of other villages' treatments given $S_{i}=s$ and $S_{i}=s^{\prime}$, respectively. As with the marginal direct effect, this definition depends on the treatment assignment mechanism.

For between-cluster indirect effects, similarly, we define the \textit{marginal between-cluster indirect effect} of $H_{i}$ on the outcome as
\begin{eqnarray*}
    \bie(h,h^{\prime}) &=& n^{-1}\sumi E_{\bA_{(-i)}\mid H_{i} = h}\left\{Y_{i}(0,S_{i},h)\right\} - n^{-1}\sumi E_{\bA_{(-i)}\mid H_{i} = h^{\prime}}\left\{Y_{i}(0,S_{i},h^{\prime})\right\},
\end{eqnarray*}
where $E_{\bA_{(-i)}\mid H_{i} = h}(\cdot)$ denotes expectation conditional on $H_{i}=h$. The quantity $\bie(h,h^{\prime})$ is the corresponding marginal effect, depending on the distribution of treatments of other villages under the condition $H_{i}=h$ and $H_{i}=h^{\prime}$.

\begin{remark}
    When the treatment assignments $A_{i}$'s are independent, e.g., if we have Bernoulli randomization within each sublocation, then the conditional expectations $E_{\bA_{(-i)}\mid S_{i} = s}(\cdot)$ and $E_{\bA_{(-i)}\mid H_{i} = h}(\cdot)$ are equivalent to expectations over the marginal distribution of $H_{i}$ and $S_{i}$, respectively.
\end{remark}

\begin{remark}
    All quantities are defined conditional on the realization of the first-stage randomization $\bZ$, which affects the distribution of $\bA$.
\end{remark}

\begin{remark}
    Here we define the two indirect effects holding the treatment status $A=0$ following \cite{hudgens2008toward}. We can similarly define analogous estimands on the treated subset.
\end{remark}

\subsection{Policy-specific causal parameters}
In random saturation designs such as \cite{egger2022general}, we can further define estimands that compare two policies, similar to \cite{hudgens2008toward}. For a specific treatment assignment policy $\psi$, define the policy-specific direct effect $\de_{\psi}$, within-cluster indirect effect $\wie_{\psi}$, and between-cluster indirect effect $\bie_{\psi}$ as follows
\begin{eqnarray*}
    \de_{\psi} &=& n^{-1}\sumi E_{\psi}\{Y_i(1,S_i,H_i)\} - n^{-1}\sumi E_{\psi}\{Y_i(0,S_i,H_i)\}, \\
    \wie_{\psi} &=& n^{-1}\sumi E_{\psi}\left\{Y_{i}(0,1,H_{i})\right\} - n^{-1}\sumi E_{\psi}\left\{Y_{i}(0,0,H_{i})\right\}, \\
    \bie_{\psi} &=& n^{-1}\sumi E_{\psi}\left\{Y_{i}(0,S_{i},1)\right\} - n^{-1}\sumi E_{\psi}\left\{Y_{i}(0,S_{i},0)\right\},
\end{eqnarray*}
where the subscript $\psi$ denotes marginalization of $(S_{i}, H_{i})$ conditional on $A_{i} = a$ under the distribution induced by policy $\psi$ for $a=0,1$ in the definition of $\de_{\psi}$, marginalization of $H_{i}$ conditional on $(A_{i},S_{i}) = (0,s)$ under the distribution induced by policy $\psi$ for $s=0,1$ in the definition of $\wie_{\psi}$, and marginalization of $S_{i}$ conditional on $(A_{i},H_{i}) = (0,h)$ under the distribution induced by policy $\psi$ for $h=0,1$ in the definition of $\bie_{\psi}$. In particular, if we take $\psi$ as the treatment policy implemented in the study, the above effects recover the in-policy causal parameters we defined in the previous section. 

For two policies $\psi_1$, $\psi_2$, we can therefore compare their direct effect and indirect effects by the contrasts $\de_{\psi_1} - \de_{\psi_2}$, $\wie_{\psi_1} - \wie_{\psi_2}$, and $\bie_{\psi_1} - \bie_{\psi_2}$, respectively.

\section{Estimation and inference}
\label{sec::estimation}

\subsection{Effect estimation}

\paragraph{Estimate village average potential outcome and controlled effects.}
We first construct an estimator for the village average potential outcome,
\begin{eqnarray*}
    \bar{Y}(a,s,h) &=& n^{-1}\sumi Y_{i}(a,s,h),
\end{eqnarray*}
for given $(a,s,h)$. Consider the Horvitz--Thompson estimator
\begin{eqnarray*}
    \hat{Y}^{\horthom}(a,s,h) &=& n^{-1}\sumi\frac{ \bbI_{i}(a,s,h)}{ \pscore{i}(a,s,h)}Y_{i},
\end{eqnarray*}
and the Hajek-type estimator
\begin{eqnarray*}
    \hat{Y}^{\haj}(a,s,h) &=& n^{-1}\sumi\frac{ \bbI_{i}(a,s,h) }{ \pscore{i}(a,s,h)}Y_{i} \Big / n^{-1}\sumi\frac{ \bbI_{i}(a,s,h)}{ \pscore{i}(a,s,h)},
\end{eqnarray*}
where $\bbI_{i}(a,s,h)=\indicator\{A_{i}=a,S_{i}=s,H_{i}=h\}$ and $\pscore{i}(a,s,h)=\pr(A_{i}=a, S_{i}=s, H_{i}=h)$ is the propensity score. These propensity scores are determined by both the experimental design and the network structure. In principle, they depend deterministically on the treatment assignment mechanism and definition of the exposure mapping, but directly calculating them can become infeasible as the treatment space grows. A practical alternative is to approximate these probabilities using Monte Carlo simulations.

We then propose the following estimators for the controlled direct and indirect effects defined in Section~\ref{sec::parameters}. For $*\in\{\horthom,\haj\}$, construct
\begin{eqnarray*}
    \hat{\de}^{*}(s,h) &=& \hat{Y}^{*}(1,s,h) - \hat{Y}^{*}(0,s,h), \\
    \hat{\wie}^{*}(s,s^{\prime},h) &=& \hat{Y}^{*}(0,s,h) - \hat{Y}^{*}(0,s^{\prime},h), \\
    \hat{\bie}^{*}(s,h,h^{\prime}) &=& \hat{Y}^{*}(0,s,h) - \hat{Y}^{*}(0,s,h^{\prime}).
\end{eqnarray*}

Under the randomization, the estimator $\hat{Y}^{\horthom}(a,s,h)$ is unbiased to $\bar{Y}(a,s,h)$. The Hajek estimator $\hat{Y}^{\haj}(a,s,h)$ is not unbiased in finite sample but is consistent to $\hat{Y}^{\haj}(a,s,h)$ and generally has more stable finite sample performance. Therefore, the corresponding Horvitz--Thompson estimators are unbiased, and the Hajek estimators are consistent.

% We also construct estimators for the aggregated versions of the village average potential outcome, aggregated at various levels. Define
% \begin{eqnarray*}
%     \hat{Y}^{\horthom}(a) &=& \frac{\sumi\indicator\{A_{i}=a\} Y_{i}}{\sumi\pr(A_{i}=a)}, \quad \hat{Y}^{\haj}(a) \ =\ \frac{\sumi\indicator\{A_{i}=a\} Y_{i}}{\sumi\indicator\{A_{i}=a\}}, \\
%     \hat{Y}^{\horthom}(a,s) &=& n^{-1}\sumi\frac{\indicator\{A_{i}=a, S_{i}=s\} }{\pr(A_{i}=a, S_{i}=s)}Y_{i}, \quad \hat{Y}^{\haj}(a,s) \ =\  n^{-1}\sumi\frac{\indicator\{A_{i}=a, S_{i}=s\} }{\pr(A_{i}=a, S_{i}=s)}Y_{i}\Big/ n^{-1}\sumi\frac{\indicator\{A_{i}=a, S_{i}=s\}}{\pr(A_{i}=a, S_{i}=s)}, \\
%     \hat{Y}^{\horthom}(a,h) &=& n^{-1}{\sumi \frac{\indicator\{A_{i}=a, H_{i}=h\}}{\pr(A_{i}=a, H_{i}=h)} }Y_{i}, \quad \hat{Y}^{\haj}(a,h) \ =\  n^{-1}{\sumi \frac{\indicator\{A_{i}=a, H_{i}=h\}}{\pr(A_{i}=a, H_{i}=h)} Y_{i}} \Big/ n^{-1}
%     {\sumi \frac{\indicator\{A_{i}=a, H_{i}=h\}}{\pr(A_{i}=a, H_{i}=h)} }.
% \end{eqnarray*}
% For $*\in\{\horthom,\haj\}$, estimators for the marginal quantities are therefore
% \begin{eqnarray*}
%     \hat{\de}^{*} &=& \hat{Y}^{*}(1) - \hat{Y}^{*}(0), \\
%     \hat{\wie}^{*}(s,s^{\prime}) &=& \hat{Y}^{*}(0,s) - \hat{Y}^{*}(0,s^{\prime}), \\
%     \hat{\bie}^{*}(h,h^{\prime}) &=& \hat{Y}^{*}(0,h) - \hat{Y}^{*}(0,h^{\prime}).
% \end{eqnarray*}
% Similarly, the corresponding Horvitz--Thompson estimators are unbiased, and the Hajek estimators are consistent.


\paragraph{Estimate in-policy and policy-specific causal effects.}
We propose to use a policy-specific re-weighting to estimate the policy-specific effects. More concretely, given weights $\gamma_i(a,s,h)$, consider a class of Horvitz--Thompson estimator with weights $\gamma_i$,
\begin{eqnarray*}
    \hat{Y}^{\horthom}(a,s,h;\Gamma) &=& n^{-1}\sumi\frac{ \bbI_{i}(a,s,h)\gamma_i(a,s,h)}{ \pscore{i}(a,s,h)}Y_{i},
\end{eqnarray*}
and the Hajek-type estimator
\begin{eqnarray*}
    \hat{Y}^{\haj}(a,s,h;\Gamma) &=& n^{-1}\sumi \gamma_i(a,s,h)\cdot n^{-1}\sumi\frac{ \bbI_{i}(a,s,h) \gamma_i(a,s,h)}{ \pscore{i}(a,s,h)}Y_{i} \Big / n^{-1}\sumi\frac{ \bbI_{i}(a,s,h) \gamma_i(a,s,h)}{ \pscore{i}(a,s,h)}. 
\end{eqnarray*}
We then propose the following estimators for the policy-specific direct and indirect effects in Section~\ref{sec::parameters}. 

Define the elements of function classes $\Gamma_{\psi}^{\de}$, $\Gamma_{\psi}^{\bie}$, and $\Gamma_{\psi}^{\wie}$ as
\begin{equation*}
    \gamma_{i,\psi}^{\de}(a,s,h) = \pscore{i,\psi}(\cdot,s,h), \quad 
    \gamma_{i,\psi}^{\wie}(a,s,h) = \pscore{i,\psi}(\cdot, \cdot, h), \quad
    \gamma_{i,\psi}^{\bie}(a,s,h) = \pscore{i,\psi}(\cdot, s, \cdot),
\end{equation*}
where $\pscore{i,\psi}(\cdot,s,h)=\pr_{\psi}(S_i=s,H_i=h)$, $\pscore{i,\psi}(\cdot, \cdot, h) = \pr_{\psi}(H_i=h)$, and $\pscore{i,\psi}(\cdot, s, \cdot)=\pr_{\psi}(S_i=s)$. For $*\in\{\horthom,\haj\}$, construct the following estimators
\begin{eqnarray*}
    \hat{\de}_{\psi}^{*} &=& \sum_{s,h=0,1} \{\hat{Y}^{*}(1,s,h;\Gamma_{\psi}^{\de})  - \hat{Y}^{*}(0,s,h;\Gamma_{\psi}^{\de})\}, \\
    \hat{\wie}_{\psi}^{*} &=& \sum_{h=0,1} \{\hat{Y}^{*}(0,1,h;\Gamma_{\psi}^{\wie}) - \hat{Y}^{*}(0,0,h;\Gamma_{\psi}^{\wie})\}, \\
    \hat{\bie}_{\psi}^{*} &=& \sum_{s=0,1} \{\hat{Y}^{*}(0,s,1;\Gamma_{\psi}^{\bie}) - \hat{Y}^{*}(0,s,0;\Gamma_{\psi}^{\bie})\}.
\end{eqnarray*}

As a special case, if we take $\phi$ as the treatment policy applied in the real study, 
for $*\in\{\horthom,\haj\}$, we have the following estimators for the in-policy marginal causal effects:
\begin{eqnarray*}
    \hat{\de}_{\phi}^{*} &=& \sum_{s,h=0,1} \{\hat{Y}^{*}(1,s,h;\Gamma_{\phi}^{\de})  - \hat{Y}^{*}(0,s,h;\Gamma_{\phi}^{\de})\}, \\
    \hat{\wie}_{\phi}^{*} &=& \sum_{h=0,1} \{\hat{Y}^{*}(0,1,h;\Gamma_{\phi}^{\wie}) - \hat{Y}^{*}(0,0,h;\Gamma_{\phi}^{\wie})\}, \\
    \hat{\bie}_{\phi}^{*} &=& \sum_{s=0,1} \{\hat{Y}^{*}(0,s,1;\Gamma_{\phi}^{\bie}) - \hat{Y}^{*}(0,s,0;\Gamma_{\phi}^{\bie})\}.
\end{eqnarray*}
\section{Theoretical properties}

\subsection{Asymptotic variance}
We now compute the asymptotic variance of the reweighted estimators. This result will be our building block for analysis. Define $\bGamma(a,s,h)=\diag\{\gamma_i(a,s,h)\}$ as the $n\times n$ diagonal matrix with diagonal terms equal to $\gamma_i(a,s,h)$ for $i=1,\ldots,n$. Meanwhile, we introduce the following $n\times n$ matrices, $\bLambda(a,s,h)$ and $\bLambda(a,s,h;\ashprime)$, that will be crucial in defining the asymptotic variance: 
\begin{eqnarray*}
    \bLambda_{(i,j)}(a,s,h) &=& \indicator\{j=i\} \frac{1 - \pscore{i}(a, s, h)}{\pscore{i}(a, s, h)} + \indicator\{j\neq i\}\frac{\pscore{ij}(a, s, h; a, s, h) - \pscore{i}(a, s, h)\pscore{j}(a, s, h)}{\pscore{i}(a, s, h)\pscore{j}(a, s, h)}, \\
    \bLambda_{(i,j)}(a,s,h;\ashprime) &=& -\indicator\{j=i\}+ \indicator\{j\neq i\}\frac{\pscore{ij}(a, s, h; \ashprime) - \pscore{i}(a, s, h)\pscore{j}(\ashprime)}{\pscore{i}(a, s, h)\pscore{j}(\ashprime)}.
\end{eqnarray*}
Also, for $*\in\{\horthom,\haj\}$, we define $Y_i^{\horthom}(a,s,h)=Y_i(a,s,h)$ as the potential outcome for unit $i$, and
\begin{eqnarray*}
    Y_i^{\haj}(a,s,h) = 
    Y_i(a,s,h) - \frac{n^{-1}\sumi \gamma_i(a,s,h)Y_i(a,s,h)}{n^{-1}\sumi\gamma_i(a,s,h)}
\end{eqnarray*}
as the centered potential outcome for unit $i$. Let $\bY^{*}(a,s,h)=(Y_1^{*}(a,s,h),\ldots,Y_n^{*}(a,s,h))^{\T}$ be the vector of potential outcomes for all units. 

\begin{theorem}\label{thm:asym_var}
For a reweighting regime $\Gamma$ and $*\in\{\horthom,\haj\}$, for a fixed treatment and exposure mapping level $(a,s,h)$, the asymptotic variance of $\hat{Y}^{*}(a,s,h;\Gamma)$ is
\begin{eqnarray*}
    && \avar\{\hat{Y}^{*}(a,s,h;\Gamma)\} \\
    &=& n^{-2}\sumi \frac{1 - \pscore{i}(a, s, h)}{\pscore{i}(a, s, h)}\{\gamma_i(a,s,h)Y_i^{*}(a, s, h)\}^2 \\
    &&+ n^{-2}\sumi\sumjneqi\frac{\pscore{ij}(a, s, h; a, s, h) - \pscore{i}(a, s, h)\pscore{j}(a, s, h)}{\pscore{i}(a, s, h)\pscore{j}(a, s, h)}\gamma_i(a,s,h)\gamma_j(a,s,h)Y_i^{*}(a, s, h)Y_j^{*}(a, s, h) \\
    &=& n^{-2} \bY^{*}(a,s,h)^{\T} \bGamma(a,s,h) \bLambda(a,s,h) \bGamma(a,s,h) \bY^{*}(a,s,h), 
\end{eqnarray*}
and for a pair $(a,s,h)$ and $(a^{\prime},s^{\prime},h^{\prime})$, the asymptotic covariance is
\begin{eqnarray*}
    && \acov\{\hat{Y}^{*}(a,s,h;\Gamma), \hat{Y}^{*}(a',s',h';\Gamma)\} \\
    &=& -n^{-2}\sumi  \gamma_i(a,s,h)\gamma_i(\ashprime) Y_i^*(a, s, h) Y_i^*(\ashprime) \\ 
    &&+ n^{-2}\sumi\sumjneqi \frac{\pscore{ij}(a, s, h; \ashprime) - \pscore{i}(a, s, h)\pscore{j}(\ashprime)}{\pscore{i}(a, s, h)\pscore{j}(\ashprime)} \gamma_i(a,s,h)\gamma_j(\ashprime)Y_i^*(a, s, h) Y_j^*(\ashprime) \\
    &=& n^{-2} \bY^*(a,s,h)^{\T} \bGamma(a,s,h)\bLambda(a,s,h;\ashprime) \bGamma(a',s',h')\bY^*(\ashprime).
\end{eqnarray*}
\end{theorem}

In practice, the second-order inclusion probabilities, $ \pi_{ij}(a,s,h;\ashprime) $, can also be simulated given a randomization policy and a network structure. See \cite{aronow2017estimating} for more discussions. 

\subsection{Consistency and asymptotic normality}

We now establish the asymptotic properties of the proposed estimators. Following the framework of \citet{aronow2017estimating} for inverse probability weighted estimators under network interference and \citet{chen2004normal} for central limit theorems under network dependence, we establish consistency and asymptotic normality under four regularity conditions.

\begin{assumption}[Bounded potential outcomes]\label{ass:bounded}
There exists a constant $C_Y < \infty$ such that $|Y_i(a,s,h)| \leq C_Y$ for all $i = 1, \ldots, n$ and all $(a,s,h) \in \{0,1\}^3$.
\end{assumption}

\begin{assumption}[Positivity of exposure probabilities]\label{ass:positivity}
There exist constants $0 < c_{\pi} < C_{\pi} < 1$ such that $c_{\pi} \leq \pscore{i}(a,s,h) \leq C_{\pi}$ for all villages $i$ and all exposure configurations $(a,s,h)$.
\end{assumption}

\begin{assumption}[Bounded network degree]\label{ass:degree}
There exists a constant $\Delta < \infty$ such that each village has at most $\Delta$ neighbors (i.e., villages within the same sublocation or within the 4 km geographic distance threshold).
\end{assumption}

\begin{assumption}[Bounded order of dependence]\label{ass:dependence}
The exposure mappings $(A_i, S_i, H_i)$ satisfy a bounded dependence condition: for any two villages $i$ and $j$ separated by graph distance greater than $m$ (where graph distance is measured on the neighborhood network), the random vectors $(A_i, S_i, H_i)$ and $(A_j, S_j, H_j)$ are independent. The constant $m < \infty$ is the order of dependence.
\end{assumption}

Under these assumptions, we establish the following asymptotic results for our policy-specific estimators.

\begin{theorem}[Consistency and asymptotic normality]\label{thm:asymptotic}
Suppose Assumptions~\ref{ass:bounded}--\ref{ass:dependence} hold. Let $\hat{\bY}^{*}_{\Gamma}$ and $\bar{\bY}^{*}_{\Gamma}$ be the vector of estimators and average weighted potential outcomes across all combinations under the weighting regime $\Gamma$, for $*\in\{\horthom, \haj\}$ Then as $n \to \infty$:
% TODO: need to modify the following statement, right now we write it as the joint convergence of the 8 estimators, so the variance should change to the variance-covariance matrices
\begin{enumerate}
    \item \textbf{Consistency:} For any weight function $\Gamma$ and any $(a,s,h)$, we have for $*\in\{\horthom, \haj\}$
    \begin{equation*}
    \hat{\bY}^{*}_\Gamma\xrightarrow{p} \bar{\bY}_\Gamma.
    \end{equation*}
    Consequently, all policy-specific contrast estimators (controlled effects, direct effects $\hat{\de}_{\psi}$, within-sublocation effects $\hat{\wie}_{\psi}$, and between-sublocation effects $\hat{\bie}_{\psi}$) are consistent for their respective population parameters.
    
    \item \textbf{Asymptotic normality:} Let $\hat{\bY}^{*}_{\Gamma}$ be the vector of estimators across all combinations, for $*\in\{\horthom, \haj\}$, 
    \begin{equation*}
    \acov(\hat{\bY}^{*}_{\Gamma})^{-1/2}\left\{\hat{\bY}^{*}_{\Gamma} - \bar{\bY}_\Gamma\right\} \xrightarrow{d} N\left(0, I_8\right).
    \end{equation*}

\end{enumerate}
\end{theorem}

\begin{proof}[Proof sketch]
The consistency result follows from the law of large numbers for weakly dependent random variables under the bounded degree and dependence assumptions. The asymptotic normality follows from the central limit theorem for $m$-dependent arrays \citep{chen2004normal}, where the dependence structure is determined by the network topology. The variance estimation consistency follows from combining the consistency of the Monte Carlo propensity score estimates (by the law of large numbers with $N = 100{,}000$ draws) and the plug-in principle for variance estimators. Detailed proofs follow the approach in \citet{aronow2017estimating}.
\end{proof}



In our empirical application with $n = 653$ villages, the bounded degree assumption is satisfied with $\Delta \approx 20$ (the maximum number of villages in any sublocation plus nearby villages within 4 km). The bounded dependence assumption holds with $m = 2$, as the exposure of village $i$ depends only on its immediate neighbors, and any village two or more steps away in the network is independent conditional on the randomization. The positivity assumption is verified empirically from the estimated propensity scores in Table~\ref{tab:propensity_scores}, where all propensity scores are bounded away from 0 and 1.


\section{Variance estimation and inference}



\subsection{Variance estimation}

The true asymptotic variances in Theorem~\ref{thm:asym_var} require knowledge of the true potential outcomes $Y_i(a,s,h)$, which are not directly observable. We propose conservative variance estimators in this section.
% need to mention the pi_i and pi_ij's are approximated using MC simulation somewhere. where?
We first construct the following variance estimator for $\var\{\hat{Y}^{*}(a,s,h;\Gamma)\}$ for $*\in\{\horthom,\haj\}$:
\begin{eqnarray*}
    \hat{\var}\{\hat{Y}^{*}(a,s,h;\Gamma)\} &=& n^{-2} (\hat{\bY}^{*})^{\T} \bGamma(a,s,h) \bOmega(a,s,h) \bGamma(a,s,h) \hat{\bY}^{*},
\end{eqnarray*}
where $\hat{\bY}^{*} = (\hat{Y}_1^{*}, \ldots, \hat{Y}_n^{*})^{\T}$, $\hat{Y}_i^{\horthom} = \bbI_{i}(a,s,h) Y_i / \pscore{i}(a,s,h)$, $\hat{Y}_i^{\haj} = \bbI_{i}(a,s,h) \tilde{Y}_i / \pscore{i}(a,s,h)$ with
$$\tilde{Y}_i = Y_i - \frac{\hat{Y}^{\haj}(a,s,h;\Gamma)}{n^{-1}\sumi \gamma_i(a,s,h)},$$
and the $n\times n$ matrix $\bOmega$ has
$$\bOmega_{(i,j)} = \frac{\pscore{ij}(a,s,h;a,s,h)-\pscore{i}(a,s,h)\pscore{j}(a,s,h)}{\pscore{ij}(a,s,h;a,s,h)}.$$

For controlled causal effects, for example, $\hat{\de}^{*}(s,h) = \hat{Y}^{*}(1,s,h) - \hat{Y}^{*}(0,s,h)$, the true asymptotic variance envolves the covariance term that depends on the joint distribution of different potential outcomes and is not identified. We employ a Cauchy--Schwarz inequality to bound the covariance term and use the following conservative variance estimator for $*\in\{\horthom,\haj\}$,
\begin{equation*}
\hat\var\{\hat{\de}^{*}(s,h)\} = \left[\hat{\textup{se}}\{\hat{Y}^{*}(1,s,h)\} + \hat{\textup{se}}\{\hat{Y}^{*}(0,s,h)\} \right]^2.
\end{equation*}
This provides a conservative variance estimator for the controlled direct effect, $\hat{\de}^{*}(s,h)$, with equality holding when the two estimators $\hat{Y}^{*}(1,s,h)$ and $\hat{Y}^{*}(0,s,h)$ are perfectly correlated. We can similarly construct conservative variance estimators for other controlled causal effects.

For policy-specific causal effects, such as the in-policy marginal direct effect $\hat{\de}^{*} = \sum_{s,h} \{\hat{Y}^{*}(1,s,h;\Gamma^{\de}_{\phi}) - \hat{Y}^{*}(0,s,h;\Gamma^{\de}_{\phi})\}$, we similarly use $\hat\var(\hat{\de}^{*}) = [\hat{\textup{se}}\{\hat{Y}^{*}(1,s,h;\Gamma^{\de}_{\phi})\} + \hat{\textup{se}}\{\hat{Y}^{*}(0,s,h;\Gamma^{\de}_{\phi})\}]^2$ as a variance estimator. We can similarly construct conservative variance estimators for other policy-specific causal effects.

\subsection{Valid inference}
In this section, we provide asymptotic results for the variance estimator to support valid inference.
\begin{corollary}[Asymptotic validity of confidence intervals]\label{cor:ci}
Under Assumptions~\ref{ass:bounded}--\ref{ass:dependence},
\begin{enumerate}
    \item The variance estimators $\hat{\var}\{\hat{Y}^{*}(a,s,h;\Gamma)\}$ for $* \in \{\horthom, \haj\}$ is consistent: 
    \begin{equation*}
    \hat{\var}\{\hat{Y}^{*}(a,s,h;\Gamma)\} \xrightarrow{p} \var\{\hat{Y}^{*}(a,s,h;\Gamma)\},
    \end{equation*}
    and therefore for a single exposure configuration $(a,s,h)$, the confidence interval
    \begin{equation*}
    \hat{Y}^{*}(a,s,h;\Gamma) \pm z_{\alpha/2} \cdot \widehat{\textup{se}}\{\hat{Y}^{*}(a,s,h;\Gamma)\}
    \end{equation*}
    achieves asymptotic coverage $1 - \alpha$ for $* \in \{\horthom, \haj\}$.
    
    \item For controlled causal effects and policy-specific causal effects, the conservative variance estimator based on the Cauchy--Schwarz bound provides asymptotically valid confidence intervals with coverage at least $1 - \alpha$. In general, the actual coverage may exceed the nominal level due to the conservativeness of the bound.
\end{enumerate}
\end{corollary}

\begin{remark}
    If the second-order inclusion probability $\pscore{ij}(a,s,h;a,s,h)$ for some pair $(i,j)$, the variance of $\hat{Y}^{*}(a,s,h)$ cannot be consistently estimated either. We can construct conservative estimators for the variance term following the estimators proposed in \cite{aronow2017estimating}.
\end{remark}


\section{Real data analysis}

\subsection{Background and experimental design}

We apply our methodology to re-analyze data  \citet{egger2022general}, which evaluated the economic impacts of a large-scale unconditional cash transfer program implemented by GiveDirectly in rural Kenya between 2014 and 2017. The study area consists of 653 villages nested within 155 sublocations across two counties in Kenya. Sublocations are administrative units, and villages within the same sublocation often share common markets, social ties, and economic connections. The treatment is an unconditional cash transfer to eligible households in treated villages.

The randomization proceeded in two stages following a saturation design. First, sublocations were randomly assigned to either high saturation (probability $p = 2/3$) or low saturation (probability $p = 1/3$). Second, within each sublocation, villages were independently assigned to treatment with probability $e_h = 2/3$ in high-saturation sublocations and $e_l = 1/3$ in low-saturation sublocations. All eligible households in treated villages received transfers. This design resulted in 328 treated villages and 325 control villages in our analysis sample.

For outcomes, we focus on four village-level enterprise outcomes measured at endline (approximately two years after the cash transfers): winsorized average profit, revenue, total cost, and wage bill acorss all enterprises in the village.

\subsection{Exposure mapping and propensity scores}

A key feature of this application is that villages may influence each other both through administrative boundaries (within the same sublocation) and through geographic proximity (nearby villages in different sublocations). To capture these two sources of interference, we construct two binary exposure variables $S_i$ and $H_i$ for each village $i$ as defined in Section~\ref{sec::parameters}. For between-sublocation spillovers, we use geographic distance data between all village pairs to construct the graph. For each village $i$, we identify its three or fewer nearest neighbors \textit{outside} its own sublocation while within distance strictly less than 4 km.

As previously mentioned, the propensity scores $\pscore{i}(a,s,h)$ are not available in closed form due to the complex dependencies between $(A_i, S_i, H_i)$ induced by the network structure. Following \citet{aronow2017estimating}, we estimate these propensity scores and the second-order inclusion probabilities $\pscore{ij}(a,s,h;\ashprime)$ via Monte Carlo simulation with $100,000$ draws. Table~\ref{tab:propensity_scores} summarizes the estimated propensity scores across the 653 villages in our sample. 

To improve computational efficiency, we exploit the sparsity structure of the network to approximate the second-order probabilities: they are only computed and stored for pairs of villages $(i,j)$ that are potentially dependent (i.e., in the same sublocation or within 4 km of each other). For all other pairs, we use the independence approximation $\pscore{ij}(a,s,h;\ashprime) \approx \pscore{i}(a,s,h) \pscore{j}(\ashprime)$. 

\begin{table}[htbp]
\centering
\caption{Summary statistics of estimated propensity scores}
\label{tab:propensity_scores}
\begin{tabular}{lccc}
\toprule
Exposure & Mean & Std & Median \\
\midrule
$(A=0,S=0,H=0)$ & 0.184 & 0.175 & 0.137 \\
$(A=0,S=0,H=1)$ & 0.143 & 0.141 & 0.099 \\
$(A=0,S=1,H=0)$ & 0.088 & 0.062 & 0.071 \\
$(A=0,S=1,H=1)$ & 0.086 & 0.068 & 0.068 \\
$(A=1,S=0,H=0)$ & 0.118 & 0.078 & 0.088 \\
$(A=1,S=0,H=1)$ & 0.098 & 0.065 & 0.078 \\
$(A=1,S=1,H=0)$ & 0.139 & 0.143 & 0.069 \\
$(A=1,S=1,H=1)$ & 0.144 & 0.154 & 0.051 \\
\midrule
\multicolumn{4}{l}{\textit{Marginal probabilities}} \\
$\pr(A_i = 1)$ & 0.499 & 0.167 & 0.500 \\
$\pr(S_i = 1)$ & 0.458 & 0.342 & 0.500 \\
$\pr(H_i = 1)$ & 0.471 & 0.224 & 0.500 \\
\bottomrule
\end{tabular}
\begin{flushleft}
\footnotesize
\textit{Notes:} Each row reports summary statistics for the propensity score $\pscore{i}(a,s,h)$ across the 653 villages. Propensity scores are estimated using 100{,}000 Monte Carlo draws. Marginal probabilities are computed by summing the relevant joint propensities for each village. All propensity scores sum to 1 within each village (not shown).
\end{flushleft}
\end{table}


\subsection{Results}

We present the results of our proposed estimators. We organize the results into two main parts: Tables~\ref{tab:controlled_effects}--\ref{tab:between_indirect} present the controlled direct and indirect effects, and Table~\ref{tab:spillover_effects} presents the in-policy marginal direct and indirect effects.

\paragraph{Controlled direct effects.}
Table~\ref{tab:controlled_effects} reports the estimated controlled direct effects $\hat{\de}^{*}(s,h)$, which measure the treatment effect of village $i$ receiving a cash transfer conditional on its exposure environment $(S_i, H_i) = (s,h)$. 

\begin{table}[H]
\centering
\caption{Controlled direct effects: $\hat{\de}^{*}(s,h)$}
\label{tab:controlled_effects}
\footnotesize
\begin{tabular}{lcccccccccccc}
\toprule
& \multicolumn{4}{c}{Horvitz--Thompson} & \multicolumn{4}{c}{Hajek} & \multicolumn{4}{c}{AIPW} \\
\cmidrule(lr){2-5} \cmidrule(lr){6-9} \cmidrule(lr){10-13}
$(S,H)$ & Profit & Revenue & Costs & Wage & Profit & Revenue & Costs & Wage & Profit & Revenue & Costs & Wage \\
\midrule
$(0,0)$ & $2{,}142$ & $3{,}739$ & $542$ & $406$ & $-360$ & $-46$ & $200$ & $163$ & $-101$ & $217$ & $172$ & $136$ \\
        & $(2{,}180)$ & $(3{,}302)$ & $(352)$ & $(262)$ & $(909)$ & $(1{,}143)$ & $(135)$ & $(116)$ & $(883)$ & $(1{,}117)$ & $(137)$ & $(117)$ \\
\addlinespace
$(0,1)$ & $3{,}211$ & $5{,}841$ & $825^{*}$ & $625^{*}$ & $-558$ & $300$ & $329^{**}$ & $268^{*}$ & $-835$ & $-124$ & $260^{*}$ & $211$ \\
        & $(2{,}406)$ & $(3{,}940)$ & $(476)$ & $(361)$ & $(616)$ & $(877)$ & $(168)$ & $(145)$ & $(559)$ & $(696)$ & $(148)$ & $(131)$ \\
\addlinespace
$(1,0)$ & $272$ & $-440$ & $-300$ & $-253$ & $3{,}539^{***}$ & $4{,}159^{***}$ & $87$ & $23$ & $2{,}817^{**}$ & $3{,}043^{***}$ & $-150$ & $-181$ \\
        & $(3{,}431)$ & $(4{,}518)$ & $(469)$ & $(372)$ & $(1{,}263)$ & $(1{,}247)$ & $(169)$ & $(159)$ & $(1{,}268)$ & $(1{,}167)$ & $(157)$ & $(149)$ \\
\addlinespace
$(1,1)$ & $-3{,}233$ & $-4{,}785$ & $-458$ & $-340$ & $-227$ & $140$ & $115$ & $73$ & $37$ & $481$ & $142$ & $95$ \\
        & $(2{,}225)$ & $(3{,}364)$ & $(353)$ & $(260)$ & $(688)$ & $(825)$ & $(158)$ & $(134)$ & $(677)$ & $(776)$ & $(153)$ & $(131)$ \\
\bottomrule
\end{tabular}
\begin{minipage}{0.95\textwidth}
\vspace{2mm}
\footnotesize
\textit{Notes:} This table reports controlled direct effects comparing treated villages ($A_i=1$) to control villages ($A_i=0$) for each exposure environment $(S_i, H_i)$, where $S_i=1$ indicates high within-sublocation treatment saturation and $H_i=1$ indicates high between-sublocation exposure to treated neighbors. All monetary values are reported in Kenyan Shillings (KES) per enterprise per month. Point estimates are shown in the first line, with robust standard errors reported in parentheses beneath. The AIPW estimator incorporates baseline covariates (enterprise counts, profit, revenue, and costs) for improved efficiency. Significance levels: $^{*}p<0.10$, $^{**}p<0.05$, $^{***}p<0.01$.

\end{minipage}
\end{table}


The results reveal several important patterns. First, comparing the three estimators, the Hajek and AIPW estimators provide substantially more precise estimates than the Horvitz--Thompson estimator, with standard errors typically 50-70\% smaller. This is consistent with the well-known finite-sample efficiency advantages of Hajek estimation and the variance reduction from covariate adjustment in AIPW \citep{aronow2017estimating}. The AIPW estimator, which incorporates baseline covariates (enterprise counts, profit, revenue, and total costs), achieves similar or slightly better precision than the Hajek estimator across most outcomes.

Second, focusing on the more precise Hajek and AIPW estimates, we find significant positive effects on profits and revenues when villages are in high within-sublocation saturation environments but low geographic exposure environments $(S=1, H=0)$: treated villages in this setting experience increases of 3,539 KES (Hajek) or 2,817 KES (AIPW) in monthly profit ($p<0.01$ for both) and 4,159 KES (Hajek) or 3,043 KES (AIPW) in monthly revenue ($p<0.01$ for both). This suggests that within-sublocation spillovers create favorable conditions for treated enterprises to thrive. 

Finally, enterprises in villages with low local saturation but high geographic exposure $(S=0,H=1)$ face higher costs and wages due to increased competition for inputs and labor, while effects are negligible or negative when both exposure levels are high or low, highlighting the need to account for multiple spillover channels.

\paragraph{Controlled indirect effects.}

We now examine the controlled indirect effects, which isolate indirect effects by holding the village's own treatment status fixed. In addition to the controlled indirect effects introduced in Section~\ref{sec::parameters}, holding treatment status $A_i=0$, we also report estimation results for analogous parameters when holding $A_i=1$. Table~\ref{tab:within_indirect} reports the within-sublocation indirect effects, which measure the impact of changing within-sublocation treatment saturation from low to high, conditional on the village's own treatment status $A_i=a$ and geographic exposure $H_i=h$.

\begin{table}[H]
\centering
\caption{Controlled within-sublocation indirect effects}
\label{tab:within_indirect}
\footnotesize
\begin{tabular}{lcccccccccccc}
\toprule
& \multicolumn{4}{c}{Horvitz--Thompson} & \multicolumn{4}{c}{Hajek} & \multicolumn{4}{c}{AIPW} \\
\cmidrule(lr){2-5} \cmidrule(lr){6-9} \cmidrule(lr){10-13}
$(A,H)$ & Profit & Revenue & Costs & Wage & Profit & Revenue & Costs & Wage & Profit & Revenue & Costs & Wage \\
\midrule
$(0,0)$ & $2{,}881$ & $4{,}420$ & $584^{*}$ & $457^{*}$ & $-182$ & $-304$ & $132$ & $134$ & $67$ & $363$ & $333^{**}$ & $304^{***}$ \\
        & $(1{,}895)$ & $(2{,}813)$ & $(337)$ & $(262)$ & $(837)$ & $(961)$ & $(170)$ & $(148)$ & $(823)$ & $(943)$ & $(159)$ & $(140)$ \\
\addlinespace
$(0,1)$ & $3{,}078$ & $4{,}848$ & $541$ & $396$ & $-1{,}136^{*}$ & $-1{,}579^{*}$ & $-92$ & $-67$ & $-1{,}333^{**}$ & $-1{,}728^{**}$ & $-76$ & $-50$ \\
        & $(2{,}072)$ & $(3{,}131)$ & $(331)$ & $(248)$ & $(671)$ & $(813)$ & $(138)$ & $(121)$ & $(654)$ & $(750)$ & $(132)$ & $(116)$ \\
\addlinespace
$(1,0)$ & $1{,}011$ & $242$ & $-258$ & $-202$ & $3{,}718^{***}$ & $3{,}902^{***}$ & $18$ & $-6$ & $2{,}985^{**}$ & $3{,}189^{**}$ & $11$ & $-13$ \\
        & $(3{,}715)$ & $(5{,}007)$ & $(483)$ & $(372)$ & $(1{,}335)$ & $(1{,}429)$ & $(134)$ & $(127)$ & $(1{,}327)$ & $(1{,}342)$ & $(134)$ & $(126)$ \\
\addlinespace
$(1,1)$ & $-3{,}366$ & $-5{,}778$ & $-741$ & $-569$ & $-805$ & $-1{,}738^{**}$ & $-306$ & $-262^{*}$ & $-461$ & $-1{,}124$ & $-194$ & $-166$ \\
        & $(2{,}558)$ & $(4{,}173)$ & $(498)$ & $(373)$ & $(633)$ & $(889)$ & $(188)$ & $(158)$ & $(581)$ & $(722)$ & $(169)$ & $(147)$ \\
\bottomrule
\end{tabular}
\begin{minipage}{0.95\textwidth}
\vspace{2mm}
\footnotesize
\textit{Notes:} This table reports controlled within-sublocation indirect effects, comparing high within-sublocation saturation ($S_i=1$) to low saturation ($S_i=0$) for each combination of own treatment status $A_i$ and geographic exposure $H_i$. All monetary values are reported in Kenyan Shillings (KES) per enterprise per month. Point estimates are shown in the first line, with robust standard errors reported in parentheses beneath. The AIPW estimator incorporates baseline covariates for improved efficiency. Significance levels: $^{*}p<0.10$, $^{**}p<0.05$, $^{***}p<0.01$.
\end{minipage}
\end{table}

The results reveal distinct patterns depending on treatment status. Among control villages ($A=0$), the Horvitz--Thompson estimator suggests large positive effects on profits and revenues from increased sublocation saturation, though with substantial uncertainty. The Hajek estimator, on the contrary, shows negative effects when geographic exposure is high $(H=1)$: control villages in high-saturation sublocations with many nearby treated villages experience profit and revenue decreases, consistent with competitive pressure from treated neighbors.

Among treated villages ($A=1$), the Hajek estimator finds large positive effects when geographic exposure is low $(H=0)$: treated villages benefit substantially from being in high-saturation sublocations when they have few treated neighbors outside their sublocation. However, when geographic exposure is also high $(H=1)$, these benefits disappear.

Table~\ref{tab:between_indirect} presents the between-sublocation indirect effects, which measure the impact of changing geographic exposure to treated villages in other sublocations from low to high, conditional on own treatment $A_i=a$ and within-sublocation saturation $S_i=s$.

\begin{table}[H]
\centering
\caption{Controlled between-sublocation indirect effects}
\label{tab:between_indirect}
\footnotesize
\begin{tabular}{lcccccccccccc}
\toprule
& \multicolumn{4}{c}{Horvitz--Thompson} & \multicolumn{4}{c}{Hajek} & \multicolumn{4}{c}{AIPW} \\
\cmidrule(lr){2-5} \cmidrule(lr){6-9} \cmidrule(lr){10-13}
$(A,S)$ & Profit & Revenue & Costs & Wage & Profit & Revenue & Costs & Wage & Profit & Revenue & Costs & Wage \\
\midrule
$(0,0)$ & $-598$ & $-964$ & $-82$ & $-52$ & $1{,}330^{**}$ & $1{,}950^{***}$ & $222^{**}$ & $176^{**}$ & $1{,}442^{***}$ & $1{,}998^{***}$ & $204^{**}$ & $161^{**}$ \\
        & $(1{,}458)$ & $(2{,}190)$ & $(233)$ & $(177)$ & $(554)$ & $(645)$ & $(93)$ & $(83)$ & $(511)$ & $(555)$ & $(86)$ & $(76)$ \\
\addlinespace
$(0,1)$ & $-402$ & $-537$ & $-125$ & $-113$ & $375$ & $674$ & $-2$ & $-25$ & $42$ & $-92$ & $-204$ & $-194$ \\
        & $(2{,}509)$ & $(3{,}755)$ & $(435)$ & $(333)$ & $(954)$ & $(1{,}128)$ & $(214)$ & $(186)$ & $(967)$ & $(1{,}138)$ & $(205)$ & $(180)$ \\
\addlinespace
$(1,0)$ & $471$ & $1{,}138$ & $200$ & $167$ & $1{,}132$ & $2{,}296^{*}$ & $350^{*}$ & $282$ & $708$ & $1{,}658$ & $293^{*}$ & $236$ \\
        & $(3{,}127)$ & $(5{,}052)$ & $(594)$ & $(446)$ & $(971)$ & $(1{,}374)$ & $(209)$ & $(178)$ & $(930)$ & $(1{,}258)$ & $(199)$ & $(172)$ \\
\addlinespace
$(1,1)$ & $-3{,}906$ & $-4{,}882$ & $-283$ & $-201$ & $-3{,}391^{***}$ & $-3{,}344^{***}$ & $26$ & $25$ & $-2{,}739^{***}$ & $-2{,}654^{***}$ & $88$ & $83$ \\
        & $(3{,}146)$ & $(4{,}127)$ & $(387)$ & $(299)$ & $(997)$ & $(944)$ & $(112)$ & $(106)$ & $(978)$ & $(805)$ & $(104)$ & $(101)$ \\
\bottomrule
\end{tabular}
\begin{minipage}{0.95\textwidth}
\vspace{2mm}
\footnotesize
\textit{Notes:} This table reports controlled between-sublocation indirect effects, comparing high geographic exposure ($H_i=1$) to low exposure ($H_i=0$) for each combination of own treatment status $A_i$ and within-sublocation saturation $S_i$. All monetary values are reported in Kenyan Shillings (KES) per enterprise per month. Point estimates are shown in the first line, with robust standard errors reported in parentheses beneath. The AIPW estimator incorporates baseline covariates for improved efficiency. Significance levels: $^{*}p<0.10$, $^{**}p<0.05$, $^{***}p<0.01$.
\end{minipage}
\end{table}

The between-sublocation indirect effects show heterogeneity across exposure configurations. Using the Hajek estimator, control villages in low-saturation sublocations $(A=0, S=0)$ experience significant positive effects from geographic exposure to treated villages. This suggests that control villages benefit from proximity to treated villages, possibly through increased economic activity, though they also face higher input costs. However, among treated villages in high-saturation sublocations $(A=1, S=1)$, the effects are strongly negative. 

\paragraph{In-policy marginal effects.} Table~\ref{tab:spillover_effects} presents estimates of the marginal direct and indirect effects using a weighting scheme induced by the implemented policy. The marginal direct effect uses weights $\gamma_i(a,s,h) = \pscore{i}(\cdot,s,h)$ to marginalize over the distribution of own treatment, and the indirect effects use weights $\pscore{i}(\cdot,s,\cdot)$ and $\pscore{i}(\cdot,\cdot,h)$. In general, the estimators for in-policy marginal effects have relatively large standard errors, making the results statistically insignificant. Again, the Hajek and AIPW estimators provide more precise estimates, with standard errors roughly 50-60\% smaller than the Horvitz--Thompson estimator. The AIPW estimator achieves the smallest standard errors overall through incorporating baseline covariates. The lack of significance also reflects the heterogeneity documented in Tables~\ref{tab:controlled_effects}--\ref{tab:between_indirect}, where the controlled causal effects demonstrate heterogeneity across exposure environments, and averaging over them reduces statistical power.

\begin{table}[H]
\centering
\caption{In-policy marginal direct and indirect effects}
\label{tab:spillover_effects}
\footnotesize
\begin{tabular}{lcccccccccccc}
\toprule
& \multicolumn{4}{c}{Horvitz--Thompson} & \multicolumn{4}{c}{Hajek} & \multicolumn{4}{c}{AIPW} \\
\cmidrule(lr){2-5} \cmidrule(lr){6-9} \cmidrule(lr){10-13}
Parameter & Profit & Revenue & Costs & Wage & Profit & Revenue & Costs & Wage & Profit & Revenue & Costs & Wage \\
\midrule
\multicolumn{13}{l}{{Marginal direct effect}} \\
 & $577$ & $1{,}180$ & $191$ & $140$ & $335$ & $810$ & $164$ & $121$ & $289$ & $655$ & $99$ & $64$ \\
 & $(2{,}597)$ & $(3{,}930)$ & $(476)$ & $(370)$ & $(1{,}051)$ & $(1{,}242)$ & $(232)$ & $(204)$ & $(994)$ & $(1{,}119)$ & $(215)$ & $(191)$ \\
\addlinespace
\multicolumn{13}{l}{{Within-sublocation indirect effect}} \\
 & $232$ & $360$ & $102$ & $89$ & $-412$ & $-639$ & $15$ & $29$ & $-378$ & $-485$ & $62$ & $67$ \\
 & $(1{,}085)$ & $(1{,}652)$ & $(200)$ & $(156)$ & $(366)$ & $(434)$ & $(88)$ & $(78)$ & $(377)$ & $(460)$ & $(83)$ & $(73)$ \\
\addlinespace
\multicolumn{13}{l}{{Between-sublocation indirect effect}} \\
 & $-374$ & $-565$ & $-49$ & $-33$ & $433$ & $663$ & $74$ & $58$ & $499$ & $682$ & $47$ & $33$ \\
 & $(1{,}385)$ & $(2{,}068)$ & $(224)$ & $(170)$ & $(548)$ & $(636)$ & $(101)$ & $(89)$ & $(514)$ & $(564)$ & $(95)$ & $(84)$ \\
\bottomrule
\end{tabular}
\begin{minipage}{0.95\textwidth}
\vspace{2mm}
\footnotesize
\textit{Notes:} This table reports in-policy marginal direct and indirect effects estimates. The marginal direct effect uses weights $\gamma_i = \pscore{i}(\cdot,s,h)$ that marginalize over the treatment distribution. Within-sublocation spillover effects compare $S=1$ to $S=0$ for control and treated villages using weights $\gamma_i = \pscore{i}(\cdot,h,\cdot)$, while between-sublocation spillover effects compare $H=1$ to $H=0$ using $\gamma_i = \pscore{i}(\cdot,\cdot,s)$. All monetary values are reported in Kenyan Shillings (KES) per enterprise per month. Point estimates are shown in the first line, with robust standard errors reported in parentheses beneath. The AIPW estimator incorporates baseline covariates for improved efficiency. Significance levels: $^{*}p<0.10$, $^{**}p<0.05$, $^{***}p<0.01$.
\end{minipage}
\end{table}

\section{Discussion}     
Propose a model-assisted approach for the estimation and inference. While we do not directly impose outcome models on our estimands, it is convenient to explore a model-assisted approach to recover the point and variance estimators \citep{zhao2022reconciling}. 

the other two were solved. maybe mention the extension to continuous exposure mapping? 

% 1. continuous exposure mapping
% 2. two bipartite networks
% 3. notation, subscript $ki$?


\bibliographystyle{apalike}
\bibliography{ref}



% \newpage 
% \appendix

% \renewcommand{\thesection}{\Alph{section}}
% \renewcommand{\thetheorem}{\Alph{section}.\arabic{theorem}}
% \setcounter{theorem}{0}
% \renewcommand{\thelemma}{\Alph{section}.\arabic{lemma}}
% \setcounter{lemma}{0}
% \renewcommand{\theproposition}{\Alph{section}.\arabic{proposition}}
% \setcounter{proposition}{0}
% \renewcommand{\thecorollary}{\Alph{section}.\arabic{corollary}}
% \setcounter{corollary}{0}
% \renewcommand{\thedefinition}{\Alph{section}.\arabic{definition}}
% \setcounter{definition}{0}
% \renewcommand{\thepage}{S\arabic{page}}
% \setcounter{page}{1}
% \renewcommand{\theequation}{\Alph{section}.\arabic{equation}}
% \setcounter{equation}{0}
% \renewcommand{\thetable}{S\arabic{table}}
% \setcounter{table}{0}

% \begin{center}
%   \LARGE {\bf Supplementary Material}
% \end{center}
\appendix 

\section{Regression-based analysis}
\subsection{Partially saturated regression with centered factors}
Run the following partially saturated regression $Y_i \sim 1+A_i + \tS_i + \tH_i + A_i\tS_i + A_i \tH_i $, with $\tS_i$ and $\tH_i$ denoting the indicators $S_i$ and $H_i$ centered at their population mean $\pi_{S,i}$ and $\pi_{H,i}$. These population means depend on the geographic network and can be computed or simulated. 

\begin{eqnarray*}
\Omega_{XX} &=& n^{-1}
\begin{pmatrix}
    n & \sumi A_i & \sumi \tS_i & \sumi \tH_i & \sumi A_i \tS_i & \sumi A_i \tH_i \\
    \sumi A_i & \sumi A_i^2 & \sumi A_i \tS_i & \sumi A_i \tH_i & \sumi A_i^2 \tS_i & \sumi A_i^2 \tH_i \\
    \sumi \tS_i & \sumi A_i \tS_i & \sumi \tS_i^2 & \sumi \tS_i \tH_i & \sumi A_i \tS_i^2 & \sumi A_i \tS_i \tH_i \\
    \sumi \tH_i & \sumi A_i \tH_i & \sumi \tS_i \tH_i & \sumi \tH_i^2 & \sumi A_i \tS_i \tH_i & \sumi A_i \tH_i^2 \\
    \sumi A_i \tS_i & \sumi A_i^2 \tS_i & \sumi A_i \tS_i^2 & \sumi A_i \tS_i \tH_i & \sumi A_i^2 \tS_i^2 & \sumi A_i^2 \tS_i \tH_i \\
    \sumi A_i \tH_i & \sumi A_i^2 \tH_i & \sumi A_i \tS_i \tH_i & \sumi A_i \tH_i^2 & \sumi A_i^2 \tS_i \tH_i & \sumi A_i^2 \tH_i^2
\end{pmatrix}
\end{eqnarray*}
Due to the demeaning step and the independence among $A_i$, $\tS_i$, and $\tH_i$, some off-diagonal values of the expectation of the $\Omega_{XX}$ matrix are zero. Define 
$\pi_A = \pr(A = 1) = pe_{h} + (1-p)e_{l}$.
The diagonal values of $\Omega_{XX}$ are:
\begin{eqnarray*}
    E\left(n^{-1} \sumi A_i^2\right) &=& \pi_A, \\
    E\left(n^{-1} \sumi \tS_i^2\right) &=& n^{-1} \sumi \pi_{S,i}(1-\pi_{S,i}), \\
    E\left(n^{-1} \sumi \tH_i^2\right) &=& n^{-1} \sumi \pi_{H,i}(1-\pi_{H,i}), \\
    E\left(n^{-1} \sumi A_i^2 \tS_i^2\right) &=& \pi_A n^{-1} \sumi \pi_{S,i}(1-\pi_{S,i}), \\
    E\left(n^{-1} \sumi A_i^2 \tH_i^2\right) &=& \pi_A n^{-1} \sumi \pi_{H,i}(1-\pi_{H,i}). 
\end{eqnarray*}
The non-zero off-diagonal terms are 
\begin{eqnarray*}
    E\left(n^{-1} \sumi A_i \tS_i^2\right) &=& \pi_A \sumi \pi_{S,i}(1-\pi_{S,i})
\end{eqnarray*}

\begin{eqnarray*}
&& E(\Omega_{XX}) \\
&=& n^{-1}
\begin{pmatrix}
    n & n\pi_A & 0 & 0 & 0 & 0 \\
    n\pi_A & n\pi_A & 0 & 0 & 0 & 0 \\
    0 & 0 & \sumi \pi_{S,i}(1-\pi_{S,i}) & 0 & \pi_A \sumi \pi_{S,i}(1-\pi_{S,i}) & 0 \\
    0 & 0 & 0 & \sumi \pi_{H,i}(1-\pi_{H,i}) & 0 & \pi_A \sumi \pi_{H,i}(1-\pi_{H,i}) \\
    0 & 0 & \pi_A \sumi \pi_{S,i}(1-\pi_{S,i}) & 0 & \pi_A\sumi \pi_{S,i}(1-\pi_{S,i}) & 0 \\
    0 & 0 & 0 & \pi_A \sumi \pi_{H,i}(1-\pi_{H,i}) & 0 & \pi_A \sumi \pi_{H,i}(1-\pi_{H,i})
\end{pmatrix}
\end{eqnarray*}

\[
E(\Omega_{XX})^{-1} =
\begin{pmatrix}
\frac{\pi_A}{n \pi_A(1 - \pi_A)} & \frac{-\pi_A}{n \pi_A(1 - \pi_A)} & 0 & 0 & 0 & 0 \\
\frac{-\pi_A}{n \pi_A(1 - \pi_A)} & \frac{1}{n \pi_A(1 - \pi_A)} & 0 & 0 & 0 & 0 \\
0 & 0 & \frac{\pi_A}{S \pi_A(1 - \pi_A)} & 0 & \frac{-\pi_A}{S \pi_A(1 - \pi_A)} & 0 \\
0 & 0 & 0 & \frac{\pi_A}{H \pi_A(1 - \pi_A)} & 0 & \frac{-\pi_A}{H \pi_A(1 - \pi_A)} \\
0 & 0 & \frac{-\pi_A}{S \pi_A(1 - \pi_A)} & 0 & \frac{1}{S \pi_A(1 - \pi_A)} & 0 \\
0 & 0 & 0 & \frac{-\pi_A}{H \pi_A(1 - \pi_A)} & 0 & \frac{1}{H \pi_A(1 - \pi_A)} \\
\end{pmatrix}.
\]




Meanwhile, we compute
\begin{eqnarray*}
    \Omega_{XY} &=& n^{-1}
\begin{pmatrix}
    \sumi Y_i & \sumi A_i Y_i & \sumi \tS_i Y_i & \sumi \tH_i Y_i & \sumi A_i \tS_i Y_i & \sumi A_i \tH_i Y_i
\end{pmatrix}^{\T}
\end{eqnarray*}
and its population mean. Taking the expectation, we have

\begin{eqnarray*}
    E\left(n^{-1}\sumi Y_i \right) &=& n^{-1}\sumi \left[\pi_A E\{Y_i(1,S_i,H_i)\}+(1-\pi_A)E\{Y_i(0,S_i,H_i)\} \right].
\end{eqnarray*}

\begin{eqnarray*}
    E\left(n^{-1}\sumi A_i Y_i \right) &=& \pi_A n^{-1}\sumi E\{Y_i(1, S_i,H_i)\}.
\end{eqnarray*}

\begin{eqnarray*}
    E\left(n^{-1}\sumi \tS_i Y_i \right) &=& \pi_A n^{-1}\sumi \pi_{S,i}(1-\pi_{S,i})E\{Y_i(1,1,H_i)\} + (1-\pi_A) n^{-1}\sumi \pi_{S,i}(1-\pi_{S,i})E\{Y_i(0,1,H_i)\} \\
    &&- \pi_A n^{-1}\sumi \pi_{S,i}(1-\pi_{S,i})E\{Y_i(1,0,H_i)\} - (1-\pi_A) n^{-1}\sumi \pi_{S,i}(1-\pi_{S,i})E\{Y_i(0,0,H_i)\} \\
    &=& \pi_A n^{-1}\sumi \pi_{S,i}(1-\pi_{S,i}) E\{Y_i(1,1,H_i) - Y_i(1,0,H_i)\} \\
    &&+ (1-\pi_A) n^{-1}\sumi \pi_{S,i}(1-\pi_{S,i}) E\{Y_i(0,1,H_i) - Y_i(0,0,H_i)\}
\end{eqnarray*}

Similarly, 
\begin{eqnarray*}
    E\left(n^{-1}\sumi \tH_i Y_i \right) &=& \pi_A n^{-1}\sumi \pi_{H,i}(1-\pi_{H,i}) E\{Y_i(1,S_i,1) - Y_i(1,S_i,0)\} \\
    &&+ (1-\pi_A) n^{-1}\sumi \pi_{H,i}(1-\pi_{H,i}) E\{Y_i(0,S_i,1) - Y_i(0,S_i,0)\}.
\end{eqnarray*}

\begin{eqnarray*}
    E\left(n^{-1}\sumi A_i \tS_i Y_i \right) &=& \pi_A n^{-1}\sumi \pi_{S,i}(1-\pi_{S,i}) E\{Y_i(1,1,H_i) - Y_i(1,0,H_i)\}, \\
    E\left(n^{-1}\sumi A_i \tH_i Y_i \right) &=& \pi_A n^{-1}\sumi \pi_{H,i}(1-\pi_{H,i}) E\{Y_i(1,S_i,1) - Y_i(1,S_i,0)\}.
\end{eqnarray*}

Define the overlap indirect effects as
\begin{eqnarray*}
    \wie^{\textup{o}} &=& \frac{n^{-1}\sumi \pi_{S,i}(1-\pi_{S,i})E\{Y_i(0,1,H_i)-Y_i(0,0,H_i)\}}{n^{-1}\sumi \pi_{S,i}(1-\pi_{S,i})}, \\
    \bie^{\textup{o}} &=& \frac{n^{-1}\sumi \pi_{H,i}(1-\pi_{H,i})E\{Y_i(0,S_i,1)-Y_i(0,S_i,0)\}}{n^{-1}\sumi \pi_{H,i}(1-\pi_{H,i})}.
\end{eqnarray*}

Therefore, the population-level values of the true coefficients from the centered regression are $(\beta_0, \beta_{A}, \beta_{S}, \beta_{H}, \beta_{AS}, \beta_{AH})^{\T}$, where
\begin{eqnarray*}
    \beta_{A} &=& \de, \\
    \beta_{S} &=& \wie^{\textup{o}}, \\
    \beta_{H} &=& \bie^{\textup{o}}, \\
    \beta_{AS} &=& \frac{\sumi \pi_{S,i}(1-\pi_{S,i}) E\{Y_i(1,1,H_i) - Y_i(1,0,H_i) - Y_i(0,1,H_i) + Y_i(0,0,H_i)\}}{\sumi \pi_{S,i}(1-\pi_{S,i})}, \\
    \beta_{AH} &=& \frac{\sumi \pi_{H,i}(1-\pi_{H,i}) E\{Y_i(1,S_i,1) - Y_i(1,S_i,0) - Y_i(0,S_i,1) + Y_i(0,S_i,0)\}}{\sumi \pi_{H,i}(1-\pi_{H,i})}
\end{eqnarray*}

To summarize, from the linear regression $Y_i \sim 1+A_i + \tS_i + \tH_i +A_i\tS_i + A_i \tH_i$, we can estimate the direct effect as the coefficient of $A_i$. For the indirect effects, because of the centering, we are not able to recover exactly $\wie$ and $\bie$, but an overlap version of the estimands.

Similarly, if we define the indirect effects based on the treated subgroup, we can recover them from $(1-\pi_A) \beta_{AS} + \beta_{S}$ and $(1-\pi_A) \beta_{AH} + \beta_{H}$.

% \subsection{Saturated regression}
% We can run 
% \begin{eqnarray*}
%     Y_i \sim 1 + A_i + S_i + H_i + A_iS_i + \ldots + A_i S_i H_i. 
% \end{eqnarray*}
% The regression is equivalent to running a stratified regression \citep{zhao2022regression}:
% \begin{eqnarray*}
%     Y_i \sim \sum_{a,s,h} \indicator\{A_i=a, S_i = s, H_i = h\}.
% \end{eqnarray*}
      

\end{document}
